% ============================================================
\section{Results}
\label{sec:results}
% ============================================================

\subsection{Test Suite}

The suite comprises 247 tests across 14 files. Module~1 tests validate compatibility scoring against known values (same key $\to 1.0$, relative major/minor $> 0.9$, tritone $< 0.5$). Module~2 tests verify beam search on synthetic graphs and mock all four ListenBrainz similarity calls. Module~3 tests cover all constraint classes with edge cases. Module~4 tests verify classifier convergence and pickle round-trips. Mocking external APIs keeps the full suite under three seconds.

\subsection{Performance}
\label{sec:results-perf}

Table~\ref{tab:benchmark} compares the Postgres mirror to the public MusicBrainz REST API.

\begin{table}[h]
\centering
\caption{Postgres mirror vs.\ public MusicBrainz REST API (mean of 3 runs).}
\label{tab:benchmark}
\begin{tabular}{@{}lrrr@{}}
\toprule
\textbf{Operation} & \textbf{API (s)} & \textbf{Postgres (s)} & \textbf{Speedup} \\
\midrule
25 single lookups & 43.5 & 24.9 & 1.75$\times$ \\
25-MBID batch & 1.16 & 0.93 & 1.25$\times$ \\
5 artist-relationship queries & 6.3 & 2.3 & 2.74$\times$ \\
\bottomrule
\end{tabular}
\end{table}

Relationship traversal sees the largest gain (2.74$\times$) because each REST call only returns one neighbourhood, so multi-hop traversal becomes a chain of round-trips, whereas the Postgres mirror executes the same traversal as a single recursive query. Typical end-to-end playlist generation runs roughly 2$\times$ faster on the mirror.

\subsection{Constraints and Mood Classification}

The min-conflicts solver achieves 100\% hard-constraint satisfaction on 50 test playlists (no repeated artists, no repeated tracks), reducing total soft-violation score by 67\% with a mean of 3.2 swap iterations to convergence. The mood classifier reached a maximum validation accuracy of 64.7\% across the six classes, with most residual confusion between perceptually adjacent moods (chill / calm). Per-model 5-fold cross-validation results are summarised in Table~\ref{tab:moodaccuracy}.

\begin{table}[h]
\centering
\caption{Mood classifier 5-fold cross-validation.}
\label{tab:moodaccuracy}
\begin{tabular}{@{}lrr@{}}
\toprule
\textbf{Model} & \textbf{Accuracy} & \textbf{Macro F1} \\
\midrule
Logistic Regression & 0.59 & 0.55 \\
MLP ($256 \to 128$) & 0.63 & 0.60 \\
Ensemble & 0.65 & 0.62 \\
\bottomrule
\end{tabular}
\end{table}

\subsection{End-to-End Generation and Demo Experience}

The demoed flow walks the user through four picker stages---genres, artists, mood journey, anchor tracks---before posting source and destination MBIDs to \texttt{/api/playlist} (\texttt{length=7}, \texttt{beam\_width=10}). Across 20 test playlists (5--15 tracks, multiple genres), every request produced a valid playlist; generation takes 3--8~s, mean 4.2~s. Beam search explored 847 transitions on average; when the local frontier ran low, ``frontier widening'' queried ListenBrainz for new candidate nodes, and 12\% of selected tracks needed Essentia fallback. The result page shows the seven-track sequence, an SVG compatibility arc, a constraint scorecard, and an auto-generated journey summary (Figure~\ref{fig:webinterface}).

\begin{figure}[ht]
\centering
\begin{tikzpicture}[
    every node/.style={font=\footnotesize, align=center},
    stage/.style={rectangle, draw, thick, fill=white, rounded corners=2pt,
                  minimum width=2.4cm, minimum height=0.8cm},
    stage-active/.style={rectangle, draw, very thick, fill=gray!22,
                          rounded corners=2pt, minimum width=2.4cm, minimum height=0.8cm},
    arr/.style={-{Stealth[length=1.6mm]}, thick, draw=black!70},
    track/.style={rectangle, draw, thin, fill=white,
                  minimum width=11cm, minimum height=0.55cm, anchor=west},
    bar/.style={rectangle, draw=none, fill=black!75, minimum height=0.18cm, anchor=west}
]

% --- 4-stage picker flow (top) ---
\node[stage]        (s1) at (-5.6, 3.3) {Genres};
\node[stage]        (s2) at (-2.8, 3.3) {Artists};
\node[stage]        (s3) at ( 0.0, 3.3) {Mood\\Journey};
\node[stage-active] (s4) at ( 2.8, 3.3) {Anchor\\Tracks};
\node[stage-active] (s5) at ( 5.6, 3.3) {Playlist};

\draw[arr] (s1) -- (s2);
\draw[arr] (s2) -- (s3);
\draw[arr] (s3) -- (s4);
\draw[arr] (s4) -- (s5);

% --- Track list (bottom) ---
\def\xL{-5.6}
\def\xR{ 5.6}
\def\rowH{0.65}

% header
\node[anchor=west] at (\xL, 1.9) {\scriptsize\textit{Source $\to$ Destination playlist (length 7, $\beta=10$)}};

% rows: y, num, title, artist, pct
\foreach \i/\y/\num/\title/\artist/\pct in {%
  1/1.4/1/Girls Just Wanna Have Fun/Cyndi Lauper/87,
  2/0.8/2/Walking on Sunshine/Katrina \& The Waves/82,
  3/0.2/3/Don't Stop Me Now/Queen/76,
  4/-0.4/4/Under Pressure/Queen \& Bowie/71,
  5/-1.0/5/Comfortably Numb/Pink Floyd/} {
    \draw[fill=gray!4, draw=black!20] (\xL, \y - 0.25) rectangle (\xR + 0.7, \y + 0.25);
    \node[anchor=west, font=\bfseries\footnotesize] at (\xL + 0.1, \y) {\num};
    \node[anchor=west, font=\footnotesize\bfseries] at (\xL + 0.7, \y) {\title};
    \node[anchor=west, font=\footnotesize\itshape, text=black!55] at (\xL + 5.4, \y) {\artist};
    \ifx\pct\empty
      \node[anchor=east, font=\scriptsize, text=black!55] at (\xR + 0.7, \y) {DEST};
    \else
      \draw[fill=black!12, draw=black!30]
        (\xR - 1.5, \y - 0.08) rectangle (\xR - 0.2, \y + 0.08);
      \pgfmathsetmacro{\barLen}{(\pct/100)*1.3}
      \draw[fill=black!75, draw=none]
        (\xR - 1.5, \y - 0.08) rectangle (\xR - 1.5 + \barLen, \y + 0.08);
      \node[anchor=west, font=\bfseries\footnotesize] at (\xR - 0.1, \y) {\pct\%};
    \fi
}

% caption-style sublabel
\node[anchor=west, font=\scriptsize\itshape, text=black!55] at (\xL, -1.7)
  {Auto-generated journey summary: \emph{``A bright pop opener winding through Queen-era arena rock into a brooding Pink Floyd close.''}};

\end{tikzpicture}
\caption{Demo result page. Top: the four-stage user flow (genres $\to$ artists $\to$ mood journey $\to$ anchor tracks) culminates in a generated playlist. Bottom: the seven-track sequence with per-transition compatibility scores (71--87\% in this run). The journey summary, full per-transition explanations, and constraint scorecard are produced alongside.}
\label{fig:webinterface}
\end{figure}
